% !TeX root = ../main.tex
\section{泥石流灾害链的动力过程模拟（问题四）}\label{sec:simulation}

\subsection{仿真区域与初始条件构建}\label{sec:simulation-domain}
% 根据 DEM、冰川地形和冰厚数据建立仿真区域，设定冰岩崩释放位置与初始规模。
\writingplaceholder

\subsection{运动控制方程与物理参数}\label{sec:governing-equations}
% 明确质量守恒、动量守恒、摩阻模型，以及必要的侵蚀卷入或水源过程。
\writingplaceholder

\subsection{数值离散与仿真方案}\label{sec:numerical-scheme}
% 说明计算网格、边界条件、时间步长、数值稳定性与各类参数的来源。
\writingplaceholder

\subsection{灾害运动过程及可视化}\label{sec:process-visualization}
% 展示不同时刻的运动范围、流深、速度及传播路径。
\writingplaceholder

\subsection{传播距离、到达时间与影响区域}\label{sec:runout-arrival}
% 定量分析主要传播过程，并将结果用于检验第五章的下游空间预警。
\writingplaceholder

\subsection{参数敏感性与情景分析}\label{sec:scenario-analysis}
% 研究初始崩塌规模、水源、摩阻、物源卷入等参数改变时的结果差异。
\writingplaceholder
